\section{Limiting resolvents and potential resampling}\label{sec:resolvents}

Fix the energy $E$ and write $\rho=\rho_{d,w}(E)>0$. We continue along
the subsequence on which the microscopic point process and marked
empirical laws of Proposition~\ref{prop:marked-limits} converge.
Proposition~\ref{prop:gaussian-marks} identifies the spectral weights at
independently sampled vertices. Our first task is to reconstruct the
resolvents from these weights, including their real additive constants.
We then use potential resampling to obtain a uniform estimate for the
scalar transform and an invariance identity for finitely many potential
changes. These are the inputs to Gaussian insertion in the next section.

The passage from Gaussian spectral weights to resolvent control, and the
use of a finite-rank change to obtain a scalar equation, follow the
methods of \cite[Section~4]{bu}. Here the weights depend on the sampled
tree environments. We establish the corresponding statements below
using the disordered local law and Gaussian-mark result already proved.

\subsection{Analytic limits and the tails of spectral sums}

Conditionally on the finite graph and its potentials, sample vertices
$v_1,v_2,\ldots$ independently and uniformly. For fixed $k$ let
$W_{n,k}$ have columns $e_{v_1},\ldots,e_{v_k}$, and define
\begin{equation}\label{eq:sampled-compression}
 M_{n,k}(z)=W_{n,k}^{\mathsf T}G(E+z/n)W_{n,k}
       =\sum_i\frac{b_i b_i^{\mathsf T}}{x_i-z},
 \qquad
 b_i=(u_i(v_1),\ldots,u_i(v_k))^{\mathsf T},\quad z\in\HH.
\end{equation}
The normalization $\|u_i\|^2=n$ cancels the factor $n$ arising from
$\lambda_i-E-z/n=(x_i-z)/n$.

These matrix-valued analytic functions are tight for uniform convergence
on compact subsets of $\HH$. Indeed,
Lemma~\ref{lem:spectral-averaging} makes $\|M_{n,k}(i)\|$ tight. If
$f(z)=\sum_j q_j/(t_j-z)$ with $q_j\ge0$, then on any fixed compact
$\mathcal K\subset\HH$ scalar kernel comparison gives
\[
 \Im f(z)\le C_{\mathcal K}\Im f(i),\qquad
 |f(z)-f(i)|\le C_{\mathcal K}\Im f(i).
\]
Apply these inequalities to real quadratic forms of $M_{n,k}$ and use
polarization. They give compact bounds, and Cauchy's estimates give
equicontinuity on smaller compact sets. Diagonal extraction therefore
allows us to retain limits $M_k$, consistently for all $k$, jointly with
the positions, all bounded spectral lists, and all fixed-radius sampled
environments. Write
\[
                         \nu=\sum_j\delta_{x_j}
\]
for the limiting microscopic point measure. This notation is distinct
from the infinite-tree density-of-states measure $\nu_{d,w}$.
An occurrence of a repeated position has its own mark throughout.

We next remove the tails in the imaginary parts and in differences of
resolvents. For $z,z_0$ in a fixed compact subset of $\HH$, and $|x|>L$
with $L$ large,
\begin{equation}\label{eq:resolvent-tail-kernels}
 \Im\frac1{x-z}
 +\left|\frac1{x-z}-\frac1{x-z_0}\right|
                  \le \frac{C}{1+x^2}.
\end{equation}
The microscopic Wegner estimate gives
$\EE\nu_n(B)\le C_w|B|$ for every bounded Borel set $B$.
Consequently, the expected scalar tail in
\eqref{eq:resolvent-tail-kernels} is $O(L^{-1})$, uniformly in $n$.
For the marked sum, use positive semidefinite residues and the identity
\[
 \EE_{v_1,\ldots,v_k}\Tr(b_i b_i^{\mathsf T})=k.
\]
Its expected norm tail is also $O_k(L^{-1})$. Truncation, followed by
convergence of bounded marked lists, proves
\begin{align}
 I(s,y)&:=\sum_j\frac{y}{(x_j-s)^2+y^2},\label{eq:limiting-poisson}\\
 \Im M_k(z)&=\sum_j b_j b_j^{\mathsf T}\Im\frac1{x_j-z},
       \label{eq:matrix-imaginary-sum}\\
 M_k(z)-M_k(z_0)
   &=\sum_j b_j b_j^{\mathsf T}
            \left(\frac1{x_j-z}-\frac1{x_j-z_0}\right).
       \label{eq:matrix-difference-sum}
\end{align}
The positive sums and the sums of absolute norms in the difference
formula are finite almost surely. The scalar imaginary trace converges
to $I(s,y)$ as well. Establishing these assertions first on countable
dense argument sets and then using the same kernel bounds makes them
simultaneous throughout $\HH$. They leave open a real constant matrix
in $M_k$; determining it requires the large-height behavior.

\subsection{Count discrepancy and scalar principal values}

The slow trace estimate in Corollary~\ref{cor:local-consequences}
passes to the limiting Poisson transform. Thus, for some $c>0$ and all
sufficiently large $y$,
\begin{equation}\label{eq:slow-limiting-trace}
 \PP\bigl(|I(s,y)-\pi\rho|>y^{-c}\bigr)\le Cy^{-c}.
\end{equation}
The constants are independent of the fixed real argument $s$.
The following elementary consequence will also fix the principal-value
normalization of the transform.

\begin{lemma}[Count discrepancy]\label{lem:discrepancy}
There is a deterministic $\chi>0$ such that, almost surely,
\begin{equation}\label{eq:discrepancy}
 \nu([v,u])=\rho(u-v)+O(R^{1-\chi}),
 \qquad -R\le v\le u\le R,\qquad R\longrightarrow\infty,
\end{equation}
uniformly in $v,u$. The implicit constant may depend on the realization.
The same conclusion holds with any endpoint convention. In particular,
\begin{equation}\label{eq:position-tail-bounds}
 \PV\sum_{|x_j|>R}\frac1{x_j}=O(R^{-\chi}),\qquad
 \sum_{|x_j|>R}\frac1{x_j^2}=O(R^{-1}).
\end{equation}
Here and below $\PV$ means symmetric truncation in the position:
$\lim_{L\to\infty}\sum_{|x_j|\le L}$.
\end{lemma}

\begin{proof}
Choose $\varepsilon>0$ sufficiently small that
$2\varepsilon<c(1-\varepsilon)$, and consider dyadic $R$.
Set $Y=R^{1-\varepsilon}$ and use a real grid of mesh
$R^{1-2\varepsilon}$ in $[-3R,3R]$. There are $O(R^{2\varepsilon})$
grid points. By \eqref{eq:slow-limiting-trace}, the sum over dyadic $R$
of their failure probabilities is finite. Moreover, changing the center
of a Poisson kernel by at most $R^{1-2\varepsilon}$ changes its ratio
by $1+O(R^{-\varepsilon})$, uniformly in the position variable.
Borel--Cantelli and positivity therefore give, for some $c_1>0$,
\begin{equation}\label{eq:poisson-grid-counts}
 I(s,Y)=\pi\rho+O(R^{-c_1}),\qquad |s|\le3R,
\end{equation}
almost surely for all sufficiently large dyadic $R$.

Every interval of length at most $Y$ in this range has $O(Y)$ points:
evaluate $I$ at its midpoint and note that each point of the interval
contributes at least a constant times $Y^{-1}$. Covering by such
intervals gives $\nu([-R,R])=O(R)$, first dyadically and then at all
large radii. It also bounds the number of points in a strip of width
$W\ge Y$, lying in $[-2R,2R]$, by $O(W)$.

For $-R\le v\le u\le R$, integrate
$\pi^{-1}I(s,Y)$ over $[v,u]$. By
\eqref{eq:poisson-grid-counts}, this gives
\[
 \int P_{v,u,Y}(x)\,d\nu(x)
     =\rho(u-v)+O(R^{1-c_1}),\qquad
 P_{v,u,Y}(x)=\frac1\pi\int_v^u
                       \frac{Y\,ds}{(x-s)^2+Y^2}.
\]
Take $W=R^{1-\varepsilon/2}$. The strips of width $W$ around the
two endpoints contain $O(W)$ points. Outside these strips, within
$[-2R,2R]$, the difference between $P_{v,u,Y}$ and the interval
indicator is $O(Y/W)$; the linear count bound makes its total
contribution $O(RY/W)$. For $|x|>2R$,
$P_{v,u,Y}(x)\le CRY/x^2$, so dyadic summation using linear counts
bounds the remaining contribution by $O(Y)$. Hence the sharp and
smoothed counts differ by $O(R^{1-\varepsilon/2})$, uniformly in
the endpoints. This also handles atoms at the endpoints. Taking
$0<\chi\le\min(c_1,\varepsilon/2)$ proves
\eqref{eq:discrepancy}, with a neighboring dyadic radius covering an
arbitrary $R$.

Apply \eqref{eq:discrepancy} separately to the positive and negative
half-lines. Their counting errors are $O(1+t^{1-\chi})$ up to distance
$t$. In the symmetric sum of $1/x_j$, the two constant-density
contributions cancel. Partial summation bounds the remaining tail by
$O(R^{-\chi})$, since the integrated error is bounded by a constant
times $\int_R^\infty t^{-1-\chi}\,dt$. The inverse-square bound follows
directly from linear counts. Possible points at zero play no role in
these tail statements.
\end{proof}

We may decrease $\chi$ whenever needed. Define
\begin{equation}\label{eq:scalar-principal-value}
                         m(z)=\PV\sum_j\frac1{x_j-z}.
\end{equation}
Lemma~\ref{lem:discrepancy} and
$1/(x-z)=1/x+O_{\mathcal K}(x^{-2})$ show that this limit exists
locally uniformly on $\HH$. It is analytic, has imaginary part
$I(\Re z,\Im z)$, and obeys
\begin{equation}\label{eq:pv-initial-height}
 m(iy)=i\pi\rho+o(1),\qquad y\longrightarrow\infty.
\end{equation}
For clarity, comparison with the transform of $\rho\,dx$ on each
half-line bounds the error by
\[
 C\int_0^\infty\frac{1+t^{1-\chi}}{t^2+y^2}\,dt
                         =O(y^{-\chi}+y^{-1}),
\]
after taking $\chi<1$. The same estimate, up to constants, holds at
$e+iy$ when $|e|\le y/2$. We will use this uniform variant to start
the height descent below. No real centering of finite scalar traces
is required: $m$ is defined directly by
\eqref{eq:scalar-principal-value}.

\subsection{The real constants in the marked resolvent}

Let $\mathcal E_\alpha$ be the full potential environment of the
$\alpha$th sampled tree, including its root potential. Write
\begin{equation}\label{eq:environment-resolvent-coefficients}
 \begin{aligned}
 m_\alpha^*=m_o(E+i0;\mathcal E_\alpha),\qquad
 &a_\alpha=\Re m_\alpha^*,\qquad
 c_\alpha=\frac{\Im m_\alpha^*}{\pi\rho},\\
 &A_k=\diag(a_1,\ldots,a_k),\quad C_k=\diag(c_1,\ldots,c_k).
 \end{aligned}
\end{equation}
These variables are finite, and $c_\alpha>0$ almost surely by
Proposition~\ref{prop:tree-stability}. The sampled environments are
iid and independent of $\nu$. Proposition~\ref{prop:gaussian-marks}
states, jointly with these data, that
\begin{equation}\label{eq:independent-gaussian-residues}
 b_j=C_k^{1/2}g_j,
 \qquad (g_j)_j\text{ independent standard Gaussian vectors in }\RR^k,
\end{equation}
conditionally on $\nu$ and the environments. Finite spectral lists,
including their multiplicities, give a consistent version of this
description on an extension of the probability space.

\begin{proposition}[Marked resolvent representation]\label{prop:marked-resolvent}
Almost surely, for every finite $k$ and every $z\in\HH$,
\begin{equation}\label{eq:marked-resolvent}
 M_k(z)=A_k+\PV\sum_j\frac{b_jb_j^{\mathsf T}}{x_j-z}.
\end{equation}
The principal values converge locally uniformly. Moreover,
$\Im M_k(z)$ is positive definite, simultaneously for all these $k,z$.
\end{proposition}

\begin{proof}
Condition on $\nu$ and the sampled environments, and put
$Z_j=b_jb_j^{\mathsf T}-C_k$. Its entries are independent across $j$,
centered, and have bounded second moments, with bounds depending on
the fixed conditioning data. By \eqref{eq:position-tail-bounds},
$\sum_j Z_j/(x_j-i)$ converges almost surely entrywise when ordered by
increasing $|x_j|$. Also
\[
                  \sum_{|x_j|>1}\frac{\|Z_j\|}{x_j^2}<\infty
\]
conditionally almost surely, because its conditional expectation is
finite. Thus the difference series between $z$ and $i$ converges
absolutely and locally uniformly. Combining it with $C_km(z)$ proves
local uniform convergence of the marked principal value, say $S_k(z)$.
Equations~\eqref{eq:matrix-imaginary-sum} and
\eqref{eq:matrix-difference-sum} show that $M_k-S_k$ is a real
symmetric matrix constant in $z$.

Linear counts imply
$\sum_j|x_j-iy|^{-2}=O(y^{-1})$. Each entry of the centered marked
sum at $iy$ consequently has conditional second moment $O(y^{-1})$.
Together with \eqref{eq:pv-initial-height}, this yields
\begin{equation}\label{eq:marked-pv-height}
                         S_k(iy)\longrightarrow i\pi\rho C_k
                         \quad\text{in probability}.
\end{equation}
The convergence remains true without conditioning, by bounded
convergence applied to the conditional probabilities.

We must compare this limit with $M_k(iy)$ on the same space as its
sampled environments. Corollary~\ref{cor:local-consequences} gives
the large-microscopic-height limit jointly with every fixed-radius
environment. Its limiting tree diagonal is the measurable function
$\mathcal E_\alpha\mapsto m_\alpha^*$ of the entire environment.
It follows that
\begin{equation}\label{eq:joint-environment-height}
 M_k(iy)\longrightarrow\diag(m_1^*,\ldots,m_k^*)
                         \quad\text{in probability}.
\end{equation}
Here is the measurable-function point in this inference. Approximate
$m_\alpha^*$ in probability, under the fixed iid environment law,
by a bounded continuous function of finitely many potential
coordinates. Joint convergence with those coordinates controls the
distance to this approximation. The two approximation errors have
the same environment law before and after the limit; letting their
probabilities tend to zero proves convergence to $m_\alpha^*$ itself.
The off-diagonal entries in \eqref{eq:joint-environment-height} tend
to zero by the sampled Ward estimate in Lemma~\ref{lem:ward}.
Thus this is a statement about the existing environments, rather than
only about the marginal distributions of the diagonals.

Subtract \eqref{eq:marked-pv-height} from
\eqref{eq:joint-environment-height}. The constant $M_k-S_k$ equals
$\diag(m_\alpha^*)-i\pi\rho C_k=A_k$ almost surely, proving
\eqref{eq:marked-resolvent}.
Finally, \eqref{eq:discrepancy} gives infinitely many positions.
Conditional on positions and environments, the first $k$ Gaussian
columns are linearly independent almost surely. Since all $c_\alpha$
are positive, their positive weighted outer products already make
$\Im M_k(z)$ positive definite for every $z\in\HH$.
\end{proof}

We have now identified the full analytic compression, including the
real part that imaginary-tail estimates could not determine. The
next argument uses this formula twice, before and after changing one
potential, to improve the scalar transform itself.

\subsection{One-site resampling and uniform transform regularity}

At the first sampled vertex, replace its potential $V_1$ by an
independent $\widehat V_1$ uniform on $[-w,w]$, and put
$D_1=\widehat V_1-V_1$. Both finite ensembles have the original
law. We may retain a joint subsequential limit of the old and new
data; a hat denotes the latter. Their sampled tree environments
differ only at the indicated root, and the finite-rank resolvent
identity gives
\begin{equation}\label{eq:rank-one-resampling}
                   (1+D_1M_1(z))\widehat M_1(z)=M_1(z).
\end{equation}

We also need a scalar comparison for this coupling. A rank-one
perturbation changes an interval count by at most two. This inequality
passes to the limiting point measures on continuity intervals, then
to arbitrary endpoint limits. Applying it to intervals between a
fixed reference point and $t$ bounds a signed cumulative count by a
constant. Partial summation against $(t-z)^{-1}$ on symmetric
truncations, with vanishing boundary terms, gives
\begin{equation}\label{eq:rank-one-scalar-difference}
             |\widehat m(s+iy)-m(s+iy)|\le \frac C y.
\end{equation}
Indeed the integral of $|(t-z)^{-2}|$ over the real line is
$\pi/y$. Both principal values exist by
Lemma~\ref{lem:discrepancy}. The same argument gives $C_k/y$ for a
change at any fixed number $k$ of sites.

\begin{proposition}[Uniform transform regularity]\label{prop:transform-regularity}
There is a deterministic $\xi>0$ such that, almost surely, for every
fixed $0<\delta<1$,
\begin{equation}\label{eq:transform-regularity}
 m(e+iy)=i\pi\rho+O(y^{-\xi}),
 \qquad |e|\le R,\quad R^\delta\le y\le R^2,
 \qquad R\longrightarrow\infty.
\end{equation}
The implicit constant may depend on $\delta$ and the realization.
\end{proposition}

\begin{proof}
We first prove a pointwise implication under a fixed bound on $m$.
For one sampled site, write
\[
 M_1(z)=a_1+c_1\bigl(m(z)+Q(z)\bigr),\qquad
 Q(z)=\PV\sum_j\frac{g_j^2-1}{x_j-z}.
\]
For every integer $p\ge1$, conditionally on $\nu$ and the environment,
\begin{equation}\label{eq:gaussian-mark-moments}
 \EE\bigl(|Q(s+iy)|^{2p}\mid\nu,\mathcal E_1\bigr)
  \le C_p\left(\sum_j|x_j-s-iy|^{-2}\right)^p
  =C_p\left(\frac{\Im m(s+iy)}y\right)^p.
\end{equation}
To verify the inequality, first use a finite sum and expand the
$2p$th absolute moment. Every surviving spectral index occurs at
least twice because $g_j^2-1$ is centered. Each partition of these
indices is bounded by the appropriate Gaussian moments and by the
$p$th power of the sum of squared weights; for a block of size
$q\ge2$, use $\sum |w_j|^q\le(\sum|w_j|^2)^{q/2}$. There are only
finitely many partitions for fixed $p$. Convergence in these moments
then proves the assertion for the series.

In particular, for every fixed $C_0,A<\infty$, at any fixed $s,y$,
\begin{equation}\label{eq:marked-concentration-implication}
 \PP\bigl(|m(s+iy)|\le C_0,\ |Q(s+iy)|>y^{-1/4}\bigr)
                              \le C_{A,C_0}y^{-A}.
\end{equation}
Choose $p\ge2A$ in \eqref{eq:gaussian-mark-moments}. The identical
assertion holds for the hatted ensemble, with any fixed enlarged
bound on $|\widehat m|$.

Choose an event $\mathcal F$, of probability $p_0>0$, on the old
environment and the independent new potential such that
\begin{equation}\label{eq:resampling-environment-event}
 |m_1^*|+|\widehat m_1^*|\le C_1,\qquad
 \Im m_1^*,\Im\widehat m_1^*\ge c_*>0,\qquad |D_1|\ge c_*.
\end{equation}
Here all constants are deterministic for the fixed positive $w$.
Such an event exists directly from the cavity formula. The branch
sum has finite entries and strictly positive imaginary part almost
surely, so with positive probability its modulus is bounded and its
imaginary part is bounded below. On that event the root Green
function is bounded, with imaginary part bounded below, uniformly
for every root potential in $[-w,w]$. Impose additionally
$|\widehat V_1-V_1|\ge w/2$, an event of positive probability.
The old environment is independent of $\nu$, and the new potential
is independent of all old data; hence $\mathcal F$ is independent
of the old point process.

Suppose $|m(z)|\le C_0$. By
\eqref{eq:rank-one-scalar-difference}, $|\widehat m(z)|$ lies in a
fixed slightly larger range for $y\ge1$. On $\mathcal F$, outside
the two exceptions in \eqref{eq:marked-concentration-implication},
insert the two marked approximations into
\eqref{eq:rank-one-resampling}. Replacing $\widehat m$ by $m$ costs
$O(y^{-1})$ and gives
\begin{equation}\label{eq:resampling-quadratic}
 \begin{split}
 P(m)&=O(y^{-1/4}),\\
 P(t)&=(1+D_1(a_1+c_1t))
                      (\widehat a_1+\widehat c_1t)-(a_1+c_1t).
 \end{split}
\end{equation}
The tree branches are unchanged by resampling, so
$(1+D_1m_1^*)\widehat m_1^*=m_1^*$. Thus the real-coefficient
polynomial $P$ vanishes at $i\pi\rho$ and at its conjugate. Its
leading coefficient is $D_1c_1\widehat c_1$, whose modulus is
bounded below on $\mathcal F$. Consequently
\[
 P(t)=D_1c_1\widehat c_1(t-i\pi\rho)(t+i\pi\rho).
\]
Since $\Im m(z)>0$, we have $|m(z)+i\pi\rho|\ge\pi\rho$.
Equation~\eqref{eq:resampling-quadratic} therefore implies
$|m(z)-i\pi\rho|\le C_2 y^{-1/4}$.

This conclusion gives an arbitrary-power failure estimate for the
old point process alone. More precisely, if
\[
 \mathcal B_z=\{|m(z)|\le C_0,
                     \ |m(z)-i\pi\rho|>C_2y^{-1/4}\},
\]
then independence gives
$\PP(\mathcal B_z\cap\mathcal F)=p_0\PP(\mathcal B_z)$.
But this intersection is contained in the union of the old and new
marked-sum exceptions, each of which has unconditional probability
at most $C_Ay^{-A}$. We have proved
\begin{equation}\label{eq:scalar-bootstrap-implication}
 \PP(\mathcal B_{s+iy})\le C_{A,C_0}y^{-A}
 \quad\text{for every }A>0,
\end{equation}
uniformly in the fixed real part $s$. In particular, this reasoning
does not require the switched process to have its original law
conditionally on the old positions or environments.

It remains to establish the bound on $m$ needed to use this
implication across an expanding region. Fix $0<\delta<1$ and use,
for dyadic $R$, a unit grid in the padded region
\[
                |e|\le2R,\qquad \tfrac12R^\delta\le y\le2R^2.
\]
There are $O(R^3)$ grid points. At its top height, the uniform
variant of \eqref{eq:pv-initial-height} gives
$m(e+iy)=i\pi\rho+o(1)$ uniformly for $|e|\le2R$.
For two heights with $y/2\le y'\le y$, the absolutely convergent
difference series gives
\begin{equation}\label{eq:height-descent-bound}
 |m(e+iy')-m(e+iy)|
       \le\frac{y-y'}{y'}\Im m(e+iy)
       \le\Im m(e+iy).
\end{equation}
Choose $C_0$ large enough that this inequality places the next
height inside $|m|\le C_0$ whenever the preceding value is within
$1$ of $i\pi\rho$. Decrease heights by one unit, shortening the
last step if necessary; for large $R$ all successive ratios are
within the range in \eqref{eq:height-descent-bound}. At each step
\eqref{eq:scalar-bootstrap-implication} then returns the value to
within $C_2 y^{-1/4}$ of $i\pi\rho$.

Choose $A$ so large that $R^{3-\delta A}$ is summable over dyadic
$R$. The union bound and Borel--Cantelli exclude all implication
failures on these grids for sufficiently large $R$, almost surely.
The preceding descent proves the estimate on the grids without
any additional trace-bound assumption. Finally,
\begin{equation}\label{eq:principal-value-derivative}
 |m'(z)|\le\sum_j|x_j-z|^{-2}=\frac{\Im m(z)}{\Im z}.
\end{equation}
Nearby positive Poisson kernels are comparable. Thus the grid
bounds on $\Im m$ and \eqref{eq:principal-value-derivative} fill
unit grid gaps at cost $O(y^{-1})$. Padding permits arbitrary
large radii. Taking, for example, any fixed $0<\xi<1/4$ absorbs
constants. Apply the argument to a countable collection of positive
$\delta$ tending to zero; for any prescribed $\delta$, a smaller
one in that collection implies the asserted estimate.
\end{proof}

\subsection{Environment-selected potential changes}

We finally record the invariance identity needed for Gaussian
insertion. Each sampled tree has a distinguished root; the
\emph{root-excluded environment} means all its potential coordinates
other than that root coordinate. A single-site map below preserves
the uniform probability measure on $[-w,w]$. The choice to apply it
may depend on the root-excluded environment, which retains the
conditional uniform law of the root potential.

\begin{proposition}[Potential-switch identity]\label{prop:potential-switch}
Fix $k$. At site $\alpha\le k$, fix a measure-preserving map
$T_\alpha:[-w,w]\to[-w,w]$ and an event $B_\alpha$ measurable in
its root-excluded tree environment. Set
\[
 D_\alpha=\boldsymbol1_{B_\alpha}
                 \bigl(T_\alpha(V_\alpha)-V_\alpha\bigr),
             \qquad D=\diag(D_1,\ldots,D_k).
\]
Then, for every finite set of arguments in $\HH$, the imaginary
parts of
\begin{equation}\label{eq:potential-switch-transform}
       m(z)-\partial_z\log\det(\Id+D M_k(z))
\end{equation}
have jointly the same law as the corresponding imaginary parts
of $m(z)$. The identity is unconditional. In particular it applies
to piecewise translations and to arbitrary measurable
root-excluded selections of whether to perform them.
\end{proposition}

\begin{proof}
First suppose the selections depend on fixed-radius potential balls
through continuity sets, and the maps are piecewise translations
with only null sets of discontinuities. At finite $n$, use these
same selections and maps when the $k$ sampled balls are disjoint
trees; otherwise make no change. Conditional on the graph, the
sampling data, and all potentials except the $k$ central ones,
every selection is fixed and every central potential remains
independent and uniform. Applying the selected maps therefore
preserves the joint finite-graph law. The probability of ball or
separation failure is $O_k(n^{-1})$ for the fixed radii.

Write $H'=H+W_{n,k}D W_{n,k}^{\mathsf T}$ on this event, and
let $x_i'$ be its microscopic eigenvalues. The determinant identity
and differentiation in the microscopic variable give exactly
\begin{equation}\label{eq:finite-potential-trace-identity}
 \sum_i\frac1{x_i'-z}
   =\sum_i\frac1{x_i-z}
                  -\partial_z\log\det(\Id+D M_{n,k}(z)).
\end{equation}
Indeed the determinant ratio is
$\det(H'-E-z/n)/\det(H-E-z/n)$; its logarithmic derivative is
$-n^{-1}\Tr G'(E+z/n)+n^{-1}\Tr G(E+z/n)$.
This also verifies the sign and the absence of any further factor
of $n$ in \eqref{eq:finite-potential-trace-identity}.

Take imaginary parts. The scalar sums converge by the tail argument
preceding \eqref{eq:limiting-poisson}; the compression functions and
their derivatives converge locally uniformly. On the coordinates
where $D_\alpha\ne0$, write the logarithmic derivative as
\begin{equation}\label{eq:switch-log-derivative}
 \Tr\bigl((D^{-1}+M_k(z))^{-1}M_k'(z)\bigr),
\end{equation}
using the corresponding compression. Positive definiteness of
$\Im M_k(z)$ makes $D^{-1}+M_k(z)$ invertible. Zero coordinates
can simply be removed. Equivalently, the determinant of
$\Id+D M_k$ is nonzero, and its logarithmic derivative has no
choice-of-branch ambiguity. The continuous mapping theorem now
passes \eqref{eq:finite-potential-trace-identity} to the limit,
jointly in any finite list of arguments, and law preservation gives
the claimed identity for these selections and maps.

For completeness, no continuity assumption on a fixed
measure-preserving map is needed in the stated result. A bounded
measurable $T_\alpha$ can, for convergence purposes, be approximated
in probability under the uniform root law by continuous functions
(or by functions with finitely many interval discontinuities).
This follows from approximation of measurable sets in Lebesgue
measure by finite unions of intervals, applied to a finite binning
of its range. The approximants need not preserve measure: the
finite-graph law-preservation identity always uses $T_\alpha$
itself. They only establish joint convergence of
$T_\alpha(V_\alpha)$ with the resolvents and environment data.
The approximation probabilities are the same at every $n$ and in
the limit because the root marginals are uniform. This proves the
fixed-radius assertion for every measurable measure-preserving map.

Finally, approximate each general $B_\alpha$ in probability by
events in finitely many root-excluded coordinates. Such
approximation follows from conditional expectations on increasing
finite-radius sigma-fields. Each finite-coordinate event can in
turn be approximated by a continuity set for its finite product
uniform law. All these approximations exclude the root potential.
For fixed $k$, the probability that any of the selected maps
differs from its approximation tends to zero. On their agreement
event, \eqref{eq:potential-switch-transform} is exactly unchanged;
therefore its finite-dimensional law converges in the approximation.
The preceding identity passes to the limit. Independent branch
orderings can be included among the sampling data when an event
uses ordered coordinates.
\end{proof}

All limits in this section use a fixed number of sampled sites.
The next section increases that number inside the already
constructed limiting law, using
Proposition~\ref{prop:potential-switch} at each finite value.
